\section{Quartic symbol of a separable dispersion}
\label{sec:h-lorentz}

Theorem~\ref{thm:ratio} is the obstruction.
Axis-aligned nearest-neighbour \(H=-A\) has \(\gamma/\beta=0\) at the band
bottom, while a rotationally invariant quartic and the relativistic series
both have \(\gamma/\beta=2\).
The quadratic piece can still be matched to \(p^2/(2m)\) with \(m=1/(4J)\).
The group of that quadratic symbol is the Schr\"odinger group.
The degree-\(4\) term breaks rotations, so the failure is not only a failure
of boosts.
On \(J_2\) the ray deflection away from the axes and the diagonals is
\(O(q^2)\), which is the same order.
A substrate whose hoppings are not axis-aligned can produce \(\gamma\neq 0\).
Those graphs are the ones that leave the square class in the interference gates.
The ratio is also what survives a quadratic rescaling between \(J_3\) and a
weave (Section~\ref{sec:h-j3w}).
